%% =====================================================================
%%  T-25-03  The Anchor
%%  -------------------------------------------------------------------
%%  One idea per file. Core is mandatory. Slots in canonical order.
%%  Max 700 words. No layout commands. No filler. New source material
%%  for this entry goes in sources/<BOOK>/T-25-03.tex -- never by
%%  rewriting this file (docs/RULES.md R0.1).
%% =====================================================================
%% @kind: topic
%% @id: T-25-03
%% @title: The Anchor
%% @subtitle: First numbers bend the reality of the price
%% @chapter: c25
%% @order: 30
%% @status: stable
%% @sources: B5
%% @updated: 2026-09-19

\topic{T-25-03}{The Anchor}{First numbers bend the reality of the price}

\begin{core}
The first number or framing on the table becomes the reference the whole bargain is measured against --- people judge outcomes as gains or losses from wherever the conversation starts, and losses loom roughly twice as large as gains. Who sets that starting point, and how, quietly decides the range in which the deal lands.
\end{core}
\begin{mechanism}
Loss aversion: the pain of losing is felt at about twice the intensity of an equivalent gain, so a counterpart shown a plausible loss will move to avoid it more eagerly than to secure a gain. Anchoring works the same ground: the first salient number drags every later judgment toward itself even when it is arbitrary. The source adds an emotional anchor --- open by naming their fears (an accusation audit), which prices the loss before any number appears --- and a semantic one: the word \emph{fair}, used to pressure (\emph{we just want a fair deal}), to accuse (\emph{that's not fair}), or to establish the legitimate frame. Compromise itself is a price: splitting the difference feels principled and lands wherever the anchors put the midpoint.
\end{mechanism}
\begin{conditions}
Strongest in price-setting, salary and terms, hostage trades --- anywhere a number or term is explicit and the counterpart cares about the reference. Weakens where the counterpart has solid market data (anchor against data), in repeated games where an extreme anchor costs reputation across rounds, and with counterparts trained to treat first numbers as artefacts.
\end{conditions}
\begin{application}
  \item Before any number, anchor the emotions: name the losses they fear, so the eventual ask reads as relief.
  \item Let them go first when you are uncertain of the range; when you do go first, go with a defensible extreme, not a rude one.
  \item Prefer a range to a point --- recall a similar deal to set the best believable ballpark, aiming the top of it at your real target.
  \item Use precise, odd figures: they signal calculation, and calculation reads as a real limit.
  \item Refuse the mid-point reflex. If their anchor is absurd, re-anchor with reasons rather than countering from it.
\end{application}
\begin{feedback}
  \item Landing: their counter lands closer to your anchor than your opener; the negotiation is happening inside your range.
  \item Landing: they start justifying against your number --- the anchor is doing the work.
  \item Failing: they ignore your number entirely and restate theirs; the anchor did not take.
  \item Failing: \emph{fair} deployed mid-negotiation with sudden moral energy --- an attempt to re-anchor the frame; answer it, never absorb it.
\end{feedback}
\begin{failure}
An indefensible extreme is read as bad faith and costs the credibility the rest of the negotiation runs on. Anchoring someone against their will in a relationship transaction converts a price talk into a trust talk. And the mirror risk: every anchor you set teaches the counterpart your real position; a counterpart who anchors you first has priced your expectations for the next three rounds.
\end{failure}
\begin{countermeasures}
  \item Refuse to negotiate inside an absurd anchor: name it, restate your basis, make them move first.
  \item Carry independent data. The only reliable antidote to a number is another number with a source.
  \item Hear \emph{fair} as a move, not a moral claim: ask exactly what fairness requires here, and who defined it.
  \item Notice the loss-framing around you --- \emph{you're about to lose this} is a price, not a fact (T-15-02).
  \item Splitting the difference is not neutrality; it is the midpoint of two anchors.
\end{countermeasures}

\sources{B5}
\seesources
\seealso{T-25-01, T-25-04, T-14-02}
